\section{地位、批评、状态表与开放问题}\label{sec:lattice-status}

本节把全章的结论汇总为一个推论、一份批评与回应、一张状态表和一份开放问题清单。

\begin{corollary}[LaMET 原理的微扰内容 \tagA]\label{cor:lamet-principle-pt}
对部分子外态、在领头幂次，\cref{sec:lamet-principle}中 LaMET 原理所断言的展开在重整化微扰论中逐阶成立，适用于：非单态夸克准 PDF（\cref{thm:lattice-lamet}）；带混合的单态夸克与胶子准 PDF（\cref{prop:lamet-gluon-singlet}）；准 GPD 与准 DA（\cref{thm:lamet-quasi-gpd}）；带约化软函数的准 TMD（\cref{thm:lamet-quasi-tmd}）。在单圈这些关系都是显式的（\cref{sec:lamet-oneloop}与\cref{thm:lamet-quasi-tmd}(b)）；在两圈，非单态 PDF 的匹配核与准 TMD 的硬核也是显式的（\cref{sec:lamet-oneloop}与\cref{thm:lamet-quasi-tmd}(c)）。（本推论只涉及部分子态。）
\uses{thm:lattice-lamet, prop:lamet-gluon-singlet, thm:lamet-quasi-gpd, thm:lamet-quasi-tmd}
\end{corollary}

\begin{proofsketch}
四类准量的陈述分别就是所引的四个结果，它们有共同的形式：准量等于微扰匹配核与光锥量的卷积（TMD 情形为乘积，并带 $\sqrt{S_r}$），加上按 $\Lambda/(xP^z)$、$M/P^z$（TMD 情形还有 $1/(b_{\mathrm T}xP^z)$）压低的幂次修正。这正是 LaMET 原理对部分子态的微扰内容。显式的单圈与两圈结果见所引各节。完整论证见所引的四个结果及其文献。
\end{proofsketch}

\paragraph{地位与缺口}
\tagA，继承四个结果中最弱的一步：共线情形的瓶颈是\cref{asm:lamet-ope}（及\cref{asm:lamet-aux-renorm}），TMD 情形是\cref{asm:lamet-order-limits}（及\cref{asm:lamet-staple}、\cref{asm:lamet-soft-rapidity}）；幂次计数各步只给出 E-S 型陈述。常说的“极限次序只改变紫外、红外物理不变”是启发性的图像；实际的论证是 OPE 路线、共线减除路线与推动／大快度路线（\cref{sec:lattice-matching}与\cref{sec:lattice-tmd}）。升级途径见\cref{rem:lamet-open}第~(1)、(2)、(7)~项。

\begin{remark}[批评与回应]\label{rem:lattice-critiques}
(a)~\emph{幂次发散的混合。}Rossi 与 Testa 指出\cite{Rossi:2017muf,Rossi:2018zkn}：在格点上 DIS 算符的高阶矩与低维算符发生幂次发散的混合，他们认为这影响等时关联函数、约化 ITD 与流—流 T 乘积三种做法。回应：非定域算符唯一的幂次发散是 Wilson 线自能，它可以因子化出去（\cite{Ji:2017rah}，据其摘要；\cref{prop:lattice-renormalizability}）；在固定 $z\neq0$ 下取连续极限后，由非定域矩阵元得到的矩是有限的，并与定域扭度 2 矩一致（\cite{Karpie:2018zaz}，据其摘要，淬火近似下的数值检验）。先在固定 $z$ 下取 $a\to0$、再在 OPE 中取 $z\to0$，这一极限次序避开了问题。这一解决方案在 2017--2018 年间有争议，双方都有正式发表的交锋：Rossi 与 Testa 称上述反驳“站不住脚”\cite[第~3.1~节]{Rossi:2018zkn}；他们随后的会议报告承认原则上可以从格点数据提取单个的矩，但坚持认为直接计算完整的 PDF 存在障碍\cite{Rossi:2018stq}（据其摘要）；回应见\cite{Karpie:2018zaz,Radyushkin:2018nbf}。据我们截至 2026 年 10 月的文献检索（INSPIRE 作者检索），未见二人此后讨论这一问题的论文。在本章的框架中争论的一部分可以说得精确：$\MSbar$ 准 PDF 的矩确实发散（\cref{prop:lamet-tails}，与\cite{Radyushkin:2018nbf}一致），但匹配定理不用这些矩，PDF 的矩由固定 $z\neq0$ 处的坐标空间因子化给出（\cref{thm:lamet-coordinate}）。

(b)~\emph{欧氏与 Minkowski。}见\cref{prop:lamet-euclid-minkowski}。

(c)~\emph{逆问题。}见\cref{rem:lamet-inverse-problem}。

(d)~\emph{格点与截面因子化。}格点结果提供的是强子矩阵元，不是对截面因子化的检验。
\end{remark}

\cref{tab:lamet-status}列出本章的全部规范陈述：带等级的陈述、数学结果与编号假设逐条列出，定义、例与不带标签的注在表末分组列出。

\begin{remark}[开放问题]\label{rem:lamet-open}
(1) 类空 Wilson 线算符的 Zimmermann 级 OPE 定理（\cref{asm:lamet-ope}）。(2) 写出 QCD 加类空静态色源的可重整性证明（\cref{asm:lamet-aux-renorm}）。(3) 扭度展开对 Ioffe 时间的一致性，哪怕只在一个可解模型中（\cref{asm:lamet-uniform}；Liu 的 Gross--Neveu 定理只覆盖虚 Ioffe 时间，第~23~章）。(4) 规范理论中强子 Wilson 线关联函数的指数衰减，且常数可用（\cref{asm:lamet-large-z}，\cref{asm:lamet-gap}）。(5) $\rho_4$ 的端点指数 $\delta$（\cref{cor:lamet-endpoint}）。(6) $x=\pm\xi$ 处的 GPD（第~23~章）。(7) 超出论证层次的准 TMD 因子化：(OL1) 超出单圈的图形证明与 (OL2) 的证明，有限长订书钉的 Glauber 分析，以及用来确定 $S_r$ 的大动量介子形状因子的因子化（第~23~章）。(8) 维数正规化 Wightman 函数的坐标空间解析性（第~23~章\cref{asm:yzliu-wightman-pt}）。(9) 更高阶：NNLO 单态与胶子核、动量空间准 GPD 的 NNLO 核、三圈重—轻核（非单态 N$^3$LO 核已有\cite{Cheng:2024wyu}，见\cref{prop:lamet-twoloop}）。(10) 带零模的扭度 3 全阶因子化。(11) 重整子处方：$m_0(\tau)$ 的普适性。(12) 对格点拟合中所用的高阶混合核检验 $\int\dd\xi\,C_{\rm hyb}=1$。(13) 动量空间幂次展开的渐近（不收敛）性质（第~23~章的 O($N$) 模型）。(14) 色非单态 DPD 与重介子 LCDA 的两步链。
\end{remark}

\begin{tcolorbox}[enhanced,breakable,title=本章澄清的常见误解,fonttitle=\bfseries,
  colback=gray!3!white,colframe=gray!55!black,boxrule=0.6pt,arc=2pt,
  left=5pt,right=5pt,top=3pt,bottom=3pt,before skip=6pt,after skip=6pt]
\textbf{“格点计算检验了截面的因子化定理。”}不对：格点提供强子矩阵元——矩、准分布、Collins--Soper 核——它们是 LaMET 匹配定理的输入与证据；截面的因子化不由格点检验（\cref{rem:lattice-critiques}(d)）。\par
\textbf{“Collins--Soper 核是微扰的。”}不对：$b_{\mathrm T}\sim1/\Lambda$ 处的 $\tilde K$ 是非微扰的，格点只能通过全阶论证级的\cref{prop:lattice-cs-kernel}（\tagA）提供它。\par
\textbf{“格点直接计算了 PDF 的 $x$ 依赖。”}不对：从有限 Ioffe 时间窗口内的数据到 $f(x)$ 要经过 Fourier 逆问题；有限窗口只确定约 $2\lambda_L/\pi$ 个模式（\cref{prop:lamet-slepian}），固定 $x$ 处的可控性依赖尾部界（\cref{prop:lamet-truncation}），而 QCD 中的尾部界是猜想（\cref{conj:lamet-decay-qcd}）。\par
\textbf{“LaMET 原理已被非微扰地证明。”}不对：\cref{conj:lamet-principle}是 \tagC；部分子态的微扰内容是 \tagA（\cref{cor:lamet-principle-pt}）。\par
\textbf{“准 PDF 的矩等于 PDF 的矩。”}不对：$\MSbar$ 准 PDF 有 $1/|x|$ 尾部，没有矩；赝 PDF 有矩，由坐标空间因子化与 PDF 的矩相联系（\cref{prop:lamet-tails}，\cref{thm:lamet-coordinate}）。\par
\textbf{“匹配核依赖强子动量 $P^z$。”}不对：核的宗量是部分子动量 $|y|P^z$（\cref{thm:lamet-parton-momentum}；以强子动量为变量的早期形式不成立，\cref{rem:lamet-early-form}）。\par
\textbf{“比值方案可以直接用于动量空间匹配。”}不对：大 $|z|$ 处比值与 RI/MOM 方案的换算因子是非微扰的（\cref{rem:lamet-ratio-large-z}，\tagC）；动量空间匹配应使用混合方案（\cref{prop:lamet-scheme-validity}）。
\end{tcolorbox}

\begin{table}[htbp]
\centering\footnotesize
\renewcommand{\arraystretch}{1.08}
\caption{第~22~章状态表。“等级”一栏对 $R_{\mathrm{cond}}$ 同时给出内容等级；“误差类型”只对最弱一步是幂次计数的陈述填写（\cref{def:intro-error-types}）；“来源”一栏：“是”表示主要来源已正式发表，“本书”表示证明在本书中完整写出，“预印本”表示主要来源只有预印本。}
\label{tab:lamet-status}
\begin{tabularx}{\textwidth}{@{}>{\raggedright\arraybackslash}p{1.45cm}>{\raggedright\arraybackslash}X>{\raggedright\arraybackslash}p{2.6cm}>{\raggedright\arraybackslash}p{3.4cm}>{\raggedright\arraybackslash}p{1.15cm}>{\raggedright\arraybackslash}p{1.0cm}@{}}
\toprule
标签 & 陈述 & 等级 & 前提／瓶颈 & 误差类型 & 来源 \\
\midrule
\cref{conj:lamet-principle} & LaMET 原理（非微扰） & \tagC & 非微扰 OPE 与因子化 & --- & 是 \\
\cref{prop:lattice-moments} & 格点可计算的矩 & \tagR\ $R_{\mathrm{OPE}}$ & $n<n_0$ 时\cref{asm:intro-growth} & --- & 是 \\
\cref{prop:lamet-support} & 支撑性质（逐图） & \tagR\ $R_{\mathrm{PT}}$ & --- & --- & 是 \\
\cref{cor:lamet-entire} & ITD 为指数型整函数 & \tagR\ $R_{\mathrm{PT}}$ & --- & --- & 是 \\
\cref{asm:lamet-aux-renorm} & QCD 加静态源可重整 & \tagA & 写出的 BRST 重整化证明 & --- & 是 \\
\cref{prop:lattice-renormalizability} & 乘法可重整性 & \tagA & \cref{asm:lamet-aux-renorm} & E-S & 是 \\
\cref{prop:lamet-ad-spacelike} & 类空与类时反常维度相同 & \tagA & \cref{asm:lamet-aux-renorm} & --- & 是 \\
\cref{lem:lamet-lorentzian} & 线性项的 Lorentz 抹平 & 数学 & --- & --- & 本书 \\
\cref{prop:lamet-lrr} & 线性重整子与 LRR & \tagP & 大 $\beta_0$；$m_0(\tau)$ 的普适性 & --- & 是 \\
\cref{prop:lamet-euclid-minkowski} & 欧氏与 Minkowski 等同 & \tagA & 作者的论证（据摘要） & --- & 是 \\
\cref{prop:lamet-oneloop} & 单圈匹配 & \tagR\ $R_{\mathrm{PT}}$ & --- & --- & 是 \\
\cref{prop:lamet-tails} & 尾部、矩与卷积定义域 & \tagR\ $R_{\mathrm{PT}}$ & --- & --- & 本书 \\
\cref{prop:lamet-twoloop} & 两圈非单态核 & \tagR\ $R_{\mathrm{PT}}$ & --- & --- & 是 \\
\cref{prop:lamet-polarized} & 极化、横度、扭度 3 & \tagP & 单圈；扭度 3 零模 & --- & 是 \\
\cref{asm:lamet-ope} & 类空 Wilson 线 OPE & \tagA & Zimmermann 级证明 & --- & 是 \\
\cref{thm:lamet-coordinate} & 坐标空间因子化 & \tagR\ $R_{\mathrm{cond}}$；内容 \tagA & \cref{asm:lamet-ope}，\cref{prop:lattice-renormalizability} & --- & 是 \\
\cref{thm:lamet-parton-momentum} & 动量空间形式 & \tagR\ $R_{\mathrm{cond}}$；内容 \tagA & 同上 & --- & 是 \\
\cref{rem:lamet-early-form} & 早期形式不成立 & \tagV & 显式单圈计算 & --- & 是 \\
\cref{prop:lamet-cfp} & Ma--Qiu 共线减除 & \tagA & 光锥规范中的减除 & E-S & 是 \\
\cref{thm:lamet-lcs} & 好的格点截面 & \tagA & \cref{thm:dis-ope} & --- & 是 \\
\cref{thm:lattice-lamet} & LaMET 匹配定理 & \tagA & \cref{asm:lamet-ope} & --- & 是 \\
\cref{prop:lamet-gluon-singlet} & 单态与胶子 & \tagA & \cref{asm:lamet-ope} & --- & 是 \\
\cref{prop:lamet-hybrid-kernel} & 混合方案单圈核 & \tagR\ $R_{\mathrm{PT}}$ & --- & --- & 是 \\
\cref{prop:lamet-scheme-validity} & 方案的适用条件 & \tagA & \cref{asm:lamet-ope} & --- & 是 \\
\cref{rem:lamet-ratio-large-z} & 比值方案大 $|z|$ 匹配 & \tagC & 非微扰的 $Z'_X$ & --- & 是 \\
\bottomrule
\end{tabularx}
\end{table}

\begin{table}[htbp]\ContinuedFloat
\centering\footnotesize
\renewcommand{\arraystretch}{1.08}
\caption[]{第~22~章状态表（续）。}
\begin{tabularx}{\textwidth}{@{}>{\raggedright\arraybackslash}p{1.45cm}>{\raggedright\arraybackslash}X>{\raggedright\arraybackslash}p{2.6cm}>{\raggedright\arraybackslash}p{3.4cm}>{\raggedright\arraybackslash}p{1.15cm}>{\raggedright\arraybackslash}p{1.0cm}@{}}
\toprule
标签 & 陈述 & 等级 & 前提／瓶颈 & 误差类型 & 来源 \\
\midrule
\cref{asm:lamet-gap} & 谱隙与边界值正则性 & \tagC & 禁闭（(GP3)） & --- & 预印本 \\
\cref{thm:lamet-exp-decay} & 指数衰减（公理化） & \tagR\ $R_{\mathrm{cond}}$；内容 \tagC & \cref{asm:intro-wightman}，\cref{asm:lamet-gap} & --- & 本书 \\
\cref{conj:lamet-decay-qcd} & QCD 中的长程形式 & \tagC & 禁闭 & --- & 预印本 \\
\cref{lem:lamet-strip} & 解析带 & 数学 & --- & --- & 本书 \\
\cref{prop:lamet-truncation} & 截断误差界 & 数学 & --- & --- & 本书 \\
\cref{prop:lamet-slepian} & 可恢复的自由度 & 数学 & --- & --- & 是 \\
\cref{prop:lamet-tmc} & 靶质量修正 & \tagR\ $R_{\mathrm{cond}}$；内容 \tagA & \cref{asm:lamet-ope} & --- & 是 \\
\cref{lem:lamet-fourier} & Fourier 引理 & 数学 & --- & --- & 本书 \\
\cref{cor:lamet-endpoint} & 端点二分 & 数学 & --- & --- & 本书 \\
\cref{prop:lamet-bvz} & 扭度 4 重整子模型 & \tagP & 大 $\beta_0$ 模型 & --- & 是 \\
\cref{asm:lamet-uniform} & Ioffe 时间一致性 & \tagC & 非微扰 OPE & --- & --- \\
\cref{asm:lamet-twist4} & 扭度 4 的结构 & \tagC & 强子扭度 4 OPE & --- & --- \\
\cref{asm:lamet-large-z} & 大距离控制 & \tagC & 定量的指数衰减 & --- & --- \\
\cref{thm:lamet-hadronic} & 强子层次匹配 & \tagR\ $R_{\mathrm{cond}}$；内容 \tagC & \cref{asm:lamet-uniform,asm:lamet-twist4,asm:lamet-large-z} & --- & 本书 \\
\cref{prop:lamet-sdf} & LaMET 与 SDF & \tagR\ $R_{\mathrm{cond}}$；内容 \tagA & \cref{asm:lamet-ope} & --- & 本书 \\
\cref{thm:lamet-quasi-gpd} & 准 GPD 与准 DA & \tagA & 非前向的\cref{asm:lamet-ope} & --- & 是 \\
\cref{prop:lamet-other-objects} & 其他准分布 & \tagP & 单圈或结构层次 & --- & 部分 \\
\cref{asm:lamet-staple} & 订书钉算符 & \tagA & 超软模式 & --- & 是 \\
\cref{asm:lamet-soft-rapidity} & 软函数的快度结构 & \tagA & 快度切片的红外步骤 & --- & 是 \\
\cref{asm:lamet-order-limits} & 极限次序 & \tagA & (OL2)；(OL1) 只验证到单圈 & --- & 是 \\
\cref{thm:lamet-quasi-tmd} & 准 TMD 因子化 & \tagA & \cref{asm:lamet-order-limits}，\cref{asm:lamet-staple}，\cref{asm:lamet-soft-rapidity} & E-S & 是 \\
\cref{prop:lattice-cs-kernel} & 格点 CS 核 & \tagA & 同上 & E-S & 是 \\
\cref{prop:lamet-tmd-ratios} & 准 TMD 之比 & \tagA & 同上 & E-S & 是 \\
\cref{cor:lamet-principle-pt} & 原理的微扰内容 & \tagA & \cref{asm:lamet-ope}，\cref{asm:lamet-order-limits} & E-S & 是 \\
\midrule
\multicolumn{6}{@{}p{\linewidth}@{}}{\emph{无标签的规范陈述。}定义：\cref{def:lattice-quasi,def:lamet-schemes,def:lamet-quasi-gpd,def:lamet-quasi-tmd,def:lamet-soft-offlc,def:lamet-reduced-soft}；显式计算的例（证明在本书中完整写出）：\cref{ex:lattice-quasi-tree,ex:lamet-hybrid-normalization,ex:lamet-fourier-numbers,ex:lamet-tmc-moments,ex:lamet-qtmd-oneloop}；说明性的注：\cref{rem:lamet-formulations,rem:lamet-domain,rem:lamet-inverse-problem,rem:lattice-critiques,rem:lamet-open}。} \\
\bottomrule
\end{tabularx}
\end{table}
